% 图 24.3 —— 按 canvas.png (1181x562) 实测重画
% 画布坐标 1pt = 1px，y 向下。
\documentclass[12pt]{standalone}
\usepackage{fontspec}
\setsansfont{TeX Gyre Heros}
\newfontfamily\mfallback{DejaVu Sans}
\usepackage{amsmath}
\usepackage{graphicx}
\usepackage{tikz}
\usetikzlibrary{arrows.meta}
\begin{document}
\begin{tikzpicture}[x=1pt, y=-1pt, line cap=round, line join=round]

  %% ================= 标签（一个标签 = 一个节点） =================
  % b x c
  \node[anchor=base west, inner sep=0pt, fill=white,
        font=\fontsize{30.8}{37}\selectfont]
    at (358,214) {\textsf{\textbf{\textit{b}}}\hspace{0.23em}{\fontsize{40}{48}\selectfont\textsf{\textbf{×}}}\hspace{0.40em}\textsf{\textbf{\textit{c}}}};
  % A (left)
  \node[anchor=base, inner sep=0pt, fill=white,
        font=\fontsize{30.8}{37}\selectfont]
    at (199.5,355) {\textsf{\textbf{\textit{A}}}};
  % b x I
  \node[anchor=base west, inner sep=0pt, fill=white,
        font=\fontsize{30.8}{37}\selectfont]
    at (385,443) {\textsf{\textbf{\textit{b}}}\hspace{0.30em}{\fontsize{40}{48}\selectfont\textsf{\textbf{×}}}\hspace{0.39em}\textsf{\textbf{\textit{I}}}};
  % A (right)
  \node[anchor=base, inner sep=0pt, fill=white,
        font=\fontsize{30.8}{37}\selectfont]
    at (876.5,313) {\textsf{\textbf{\textit{A}}}};
  % pi_1(A) x 0  (left)
  \node[anchor=base west, inner sep=0pt, fill=white,
        font=\fontsize{30.8}{37}\selectfont]
    at (130,542) {\scalebox{0.72}[1]{\mfallback{\textbf{\textit{π}}}}\raisebox{-11.5pt}{\fontsize{24}{28}\selectfont\textsf{\textbf{1}}}\hspace{-0.06em}\textsf{\textbf{(}}\textsf{\textbf{\textit{A}}}\textsf{\textbf{)}}\hspace{0.14em}{\fontsize{40}{48}\selectfont\textsf{\textbf{×}}}\hspace{0.24em}\textsf{\textbf{0}}};
  % pi_1(A) x 0  (right)
  \node[anchor=base west, inner sep=0pt, fill=white,
        font=\fontsize{30.8}{37}\selectfont]
    at (801,542) {\scalebox{0.72}[1]{\mfallback{\textbf{\textit{π}}}}\raisebox{-11.5pt}{\fontsize{24}{28}\selectfont\textsf{\textbf{1}}}\hspace{-0.13em}\textsf{\textbf{(}}\textsf{\textbf{\textit{A}}}\textsf{\textbf{)}}\hspace{0.20em}{\fontsize{40}{48}\selectfont\textsf{\textbf{×}}}\hspace{0.20em}\textsf{\textbf{0}}};
  % b x 0
  \node[anchor=base west, inner sep=0pt, fill=white,
        font=\fontsize{30.8}{37}\selectfont]
    at (1011,487) {\textsf{\textbf{\textit{b}}}\hspace{0.21em}{\fontsize{40}{48}\selectfont\scalebox{1.45}[1]{\textsf{\textbf{×}}}}\hspace{0.12em}\textsf{\textbf{0}}};

  %% ================= 框与轴 =================
  \draw[line width=4.6pt] (17.5,16.5) rectangle (508.5,502.5);
  \draw[line width=4.6pt] (673,16.5) rectangle (1167,502.5);
  \draw[line width=4.6pt] (340.5,16.5) -- (340.5,525);
  \draw[line width=4.6pt] (1001.5,16.5) -- (1001.5,502.5);

  %% ================= 左图 A：四边形 + 斜线 =================
  \def\Apath{(124.4,157.6) -- (340.5,201.9) -- (340.5,282.9) -- (67.6,352.1) -- cycle}
  \begin{scope}
    \clip \Apath;
    \foreach \k in {-4,...,16} {
      \pgfmathsetmacro{\c}{344.5+18.35*\k}
      \draw[line width=2.9pt] (\c,0) -- (0,\c);
    }
  \end{scope}
  \draw[line width=5.0pt] \Apath;

  %% ================= 右图 A：自由闭合曲线 + 斜线 =================
  \def\blobpath{plot[smooth cycle] coordinates {
    (759.4,151.0) (765.8,145.0) (774.2,142.1) (783.1,141.8) (791.6,144.5)
    (799.4,149.0) (806.9,153.8) (814.7,158.2) (822.7,162.3) (831.0,165.7)
    (839.6,168.3) (848.4,169.7) (857.3,169.3) (865.9,166.9) (873.9,162.8)
    (881.3,157.8) (888.2,152.1) (894.8,146.1) (901.3,139.9) (907.9,133.8)
    (914.7,128.0) (921.9,122.7) (929.9,118.5) (938.5,116.1) (947.4,116.1)
    (956.1,118.3) (964.2,122.0) (971.7,126.8) (978.6,132.6) (984.5,139.3)
    (990.1,147.0) (995.0,155.4) (996.3,164.3) (996.5,173.2) (996.3,182.2)
    (996.1,191.1) (995.3,200.1) (995.9,209.0) (996.4,218.0) (996.7,226.9)
    (996.2,235.9) (994.6,244.7) (991.0,252.9) (985.3,259.7) (976.2,265.1)
    (968.1,269.1) (959.6,273.5) (950.9,275.7) (942.0,277.2) (933.0,274.7)
    (924.1,274.5) (915.1,274.2) (906.1,274.0) (897.2,274.3) (888.3,275.3)
    (879.5,277.0) (870.9,279.4) (862.5,282.5) (854.3,286.2) (846.5,290.6)
    (839.1,295.6) (832.1,301.2) (825.5,307.2) (819.1,313.5) (812.9,320.0)
    (806.8,326.5) (800.6,333.0) (794.1,339.2) (786.6,343.9) (777.9,345.8)
    (769.0,344.7) (760.6,341.6) (753.1,336.7) (746.8,330.4) (741.7,323.1)
    (738.0,314.9) (735.7,306.3) (735.0,297.3) (735.8,288.4) (737.9,279.7)
    (741.1,271.4) (745.0,263.3) (749.3,255.5) (754.0,247.8) (758.8,240.2)
    (763.5,232.6) (767.9,224.8) (771.8,216.8) (774.8,208.4) (776.6,199.6)
    (776.8,190.6) (775.0,181.9) (771.2,173.8) (765.5,166.9) (760.3,159.7)
  }}
  \begin{scope}
    \clip \blobpath;
    \foreach \k in {-4,...,14} {
      \pgfmathsetmacro{\c}{978.2+18.2*\k}
      \draw[line width=2.9pt] (\c,0) -- (0,\c);
    }
  \end{scope}
  \draw[line width=5.6pt] \blobpath;

  %% ================= 点、端点记号、箭头 =================
  \fill (340.5,202) circle[radius=7.2];
  \fill (1001.5,502) circle[radius=7.2];

  % 左图闭区间 [ ]
  \draw[line width=5pt] (83,482.5) -- (68.5,482.5) -- (68.5,523.5) -- (82,523.5);
  \draw[line width=5pt] (324,483) -- (341,483) (324,522.5) -- (341,522.5);

  % 右图开区间 ( )
  \draw[line width=5pt] (743.5,483.5) arc[start angle=237.6, end angle=121.4, radius=23.7];
  \draw[line width=5pt] (981,483) arc[start angle=-77.5, end angle=77.5, radius=21.1];

  % b x I 的箭头
  \draw[-{Stealth[length=35pt,width=24pt]}, line width=3.4pt]
    (384,433) .. controls (370,417) and (356,396) .. (346,365);

  % 钉死画布
  \pgfresetboundingbox
  \path[use as bounding box] (0,0) rectangle (1181,562);
\end{tikzpicture}
\end{document}
